

\begin{figure*}[t]
\centering
\footnotesize

\begin{tcolorbox}[
    boxsep=0.5pt,
    left=1pt,right=1pt,top=1pt,bottom=1pt,
    colback=blue!5!white,
    colframe=blue!40!black,
    before skip=1pt,
    after skip=1pt,
    sharp corners=downhill
]
\textbf{Question One:} What is the IFRS 16 treatment of short term or low value leased assets? What are the restrictions?
\end{tcolorbox}

\begin{tcolorbox}[
    boxsep=0.5pt,
    left=1pt,right=1pt,top=1pt,bottom=1pt,
    colback=green!2!white,
    colframe=green!25!black,
    before skip=1pt,
    after skip=1pt,
    sharp corners=downhill
]
\textbf{Gold Response:} IFRS 16 Leases permit a simplified treatment for assets with a lease period of 12 months or less, or of low value. Although the standard does not give a numerical definition of \dots
\end{tcolorbox}

\begin{tcolorbox}[
    boxsep=0.5pt,
    left=1pt,right=1pt,top=1pt,bottom=1pt,
    colback=blue!5!white,
    colframe=blue!40!black,
    before skip=1pt,
    after skip=1pt,
    sharp corners=downhill,
]
\textbf{Question Two:} On April 1st 2023, Abigail acquired telephones under a two-year lease agreement. The terms of the lease require an initial payment of \$3,000, followed by a payment of \$9,000 on March 31st 2024 and an additional \$9,000 on March 31st 2025. Show the impact of this lease arrangement on Abigail's financial statements for the year ended 31 December 2023 under IFRS 16?
\end{tcolorbox}

\begin{tcolorbox}[
    boxsep=0.5pt,
    left=1pt,right=1pt,top=1pt,bottom=1pt,
    colback=green!2!white,
    colframe=green!25!black,
    before skip=1pt,
    after skip=1pt,
    sharp corners=downhill,
]
\textbf{Gold Response:}
Annual lease rental expense $= (3000 + 9000 + 9000)/2 = 10500$ per annum.
Expense to 31 December 2023 $= 10500 \times \frac{9}{12} = 7875$.
Accrued expense $= 7875 - 3000 = 4875$.

The expense in this period of \$7,875 is not the same as the payment of \$3,000, so we need to accrue an additional expense of \$4,875.
\end{tcolorbox}

\caption{Truncated example of a question from BFF-Bench. For each question, we include a human-written gold answer that contains the final answer and a complete reasoning chain. The full version of this example is shown in \cref{fig:bffbench-sample-full} in \cref{sec:bff-bench-examples}.}
\label{fig:bffbench-sample}
\end{figure*}
